\documentclass[11pt,a4paper]{article}

\usepackage[utf8]{inputenc}
\usepackage[T1]{fontenc}
\usepackage[ngerman]{babel}
\usepackage{amsmath, amssymb, amsthm}
\usepackage{booktabs}
\usepackage{enumitem}
\usepackage{geometry}
\usepackage{hyperref}

\geometry{margin=2.5cm}
\setlist[itemize]{topsep=4pt,itemsep=2pt}
\setlist[enumerate]{topsep=4pt,itemsep=3pt}

\newtheorem{definition}{Definition}

\title{Thema 5: Time Series Cross-Validation for Temporal Data}
\author{Projektgruppe Statistical Machine Learning}
\date{\today}

\begin{document}
\maketitle

\section{Projektziel}

Ziel der Projektarbeit ist es, verschiedene Cross-Validation-Verfahren fuer
Zeitreihendaten systematisch zu vergleichen. Im Mittelpunkt steht die Frage,
welches Verfahren bei zeitlich abhaengigen Daten die zuverlaessigsten
Fehlerschaetzungen und Modellentscheidungen liefert.

\section{Forschungsfrage}

\begin{quote}
Welche Cross-Validation-Methode eignet sich fuer Forecasting-Aufgaben mit
zeitlicher Abhaengigkeit am besten, wenn Modelle anhand ihres geschätzten
Vorhersagefehlers ausgewaehlt werden?
\end{quote}

Verglichen werden klassische und zeitreihenspezifische Verfahren:
\begin{itemize}
  \item k-fold Cross-Validation,
  \item Leave-One-Out Cross-Validation (LOOCV),
  \item Rolling-Origin Cross-Validation,
  \item Blocked Cross-Validation,
  \item h-block Cross-Validation.
\end{itemize}

\section{Wichtige Definitionen}

\begin{definition}[Zeitreihe]
Eine Zeitreihe ist eine geordnete Folge von Beobachtungen
\[
  y_1, y_2, \ldots, y_T,
\]
bei der die Reihenfolge der Beobachtungen inhaltlich relevant ist. Anders als
bei unabhaengigen Querschnittsdaten kann $y_t$ von vorherigen Werten
$y_{t-1}, y_{t-2}, \ldots$ abhaengen.
\end{definition}

\begin{definition}[Forecasting-Modell]
Ein Forecasting-Modell verwendet Informationen aus der Vergangenheit, um einen
zukuenftigen Wert zu prognostizieren. Ein einfaches Lag-Modell mit $p$ Lags hat
die Form
\[
  y_t = \beta_0 + \beta_1 y_{t-1} + \cdots + \beta_p y_{t-p} + \varepsilon_t.
\]
\end{definition}

\begin{definition}[Train-Test-Split fuer Zeitreihen]
Bei Zeitreihen muss die zeitliche Ordnung erhalten bleiben. Ein typischer Split
nutzt die ersten $T_{\text{train}}$ Beobachtungen als Trainingsdaten und die
letzten $T_{\text{test}}$ Beobachtungen als Hold-out-Testdaten:
\[
  \{1,\ldots,T_{\text{train}}\}
  \quad \text{und} \quad
  \{T_{\text{train}}+1,\ldots,T\}.
\]
\end{definition}

\begin{definition}[Cross-Validation-Fehler]
Der Cross-Validation-Fehler ist der durchschnittliche Vorhersagefehler ueber
mehrere Validierungssplits:
\[
  \widehat{Err}_{CV}
  =
  \frac{1}{K}\sum_{k=1}^{K}
  \frac{1}{|V_k|}
  \sum_{i \in V_k}
  L(y_i, \hat{f}^{(-k)}(x_i)),
\]
wobei $V_k$ das Validierungsset im Split $k$ ist und $L(\cdot)$ eine
Verlustfunktion, z.B. quadratischer Fehler.
\end{definition}

\begin{definition}[k-fold Cross-Validation]
Bei k-fold CV werden die Daten in $K$ Folds aufgeteilt. Jeder Fold wird einmal
als Validierungsset genutzt, die restlichen Folds als Trainingsset. Bei
Zeitreihen ist eine zufaellige Fold-Bildung problematisch, weil dadurch
Information aus der Zukunft indirekt in das Training gelangen kann.
\end{definition}

\begin{definition}[LOOCV]
Leave-One-Out Cross-Validation ist der Spezialfall von k-fold CV mit
$K = n$. Pro Split wird genau eine Beobachtung validiert. Fuer Zeitreihen ist
LOOCV kritisch, weil benachbarte Zeitpunkte stark abhaengig sein koennen.
\end{definition}

\begin{definition}[Rolling-Origin Cross-Validation]
Rolling-Origin CV respektiert die Zeitordnung. Das Trainingsfenster besteht aus
vergangenen Beobachtungen, das Validierungsfenster liegt zeitlich danach:
\[
  \{1,\ldots,t\} \rightarrow \{t+1,\ldots,t+h\}.
\]
Das Trainingsfenster kann expandierend oder rollierend gewaehlt werden.
\end{definition}

\begin{definition}[Blocked Cross-Validation]
Blocked CV teilt die Zeitreihe in zusammenhaengende Bloecke. Ganze Bloecke
werden als Validierungsdaten verwendet. Dadurch bleibt lokale zeitliche Struktur
besser erhalten als bei zufaelligen Folds.
\end{definition}

\begin{definition}[h-block Cross-Validation]
Bei h-block CV wird zwischen Trainings- und Validierungsdaten eine Luecke der
Laenge $h$ gelassen. Diese Luecke soll die Abhaengigkeit zwischen Training und
Validierung reduzieren.
\end{definition}

\section{Daten-generierende Prozesse}

Da das Ziel ein methodischer Vergleich ist, werden Zeitreihen simuliert. Dadurch
ist der wahre Prozess bekannt und die Verfahren koennen unter kontrollierten
Bedingungen verglichen werden.

\subsection{AR(1)-Prozess}
\[
  y_t = \phi y_{t-1} + \varepsilon_t,
  \qquad
  \varepsilon_t \sim \mathcal{N}(0,\sigma^2).
\]
Der Parameter $\phi$ steuert die zeitliche Abhaengigkeit.

\subsection{MA(q)-Prozess}
\[
  y_t = \varepsilon_t + \theta_1\varepsilon_{t-1}
  + \cdots + \theta_q\varepsilon_{t-q}.
\]
Der aktuelle Wert haengt von aktuellen und vergangenen Schocks ab.

\subsection{ARMA(p,q)-Prozess}
\[
  y_t =
  \phi_1 y_{t-1} + \cdots + \phi_p y_{t-p}
  + \varepsilon_t
  + \theta_1\varepsilon_{t-1}
  + \cdots + \theta_q\varepsilon_{t-q}.
\]
ARMA-Prozesse kombinieren autoregressive und Moving-Average-Komponenten.

\subsection{Trend und Saisonalitaet}
Ein Prozess mit Trend und saisonaler Komponente kann z.B. so simuliert werden:
\[
  y_t = \alpha + \gamma t + A \sin\left(\frac{2\pi t}{s}\right)
  + \varepsilon_t.
\]
Dabei beschreibt $\gamma$ den Trend, $A$ die saisonale Amplitude und $s$ die
Periodenlaenge.

\section{Kandidatenmodelle}

In jeder Simulation werden mehrere Forecasting-Modelle trainiert und verglichen.
Die Modelle sollen unterschiedlich flexibel sein:

\begin{itemize}
  \item lineare Lag-Regression,
  \item Ridge Regression mit Lags,
  \item Lasso Regression mit Lags,
  \item Regression Tree,
  \item Random Forest oder ein anderes tree-based Modell.
\end{itemize}

Als Prädiktoren werden vergangene Werte der Zielvariable verwendet:
\[
  x_t = (y_{t-1}, y_{t-2}, \ldots, y_{t-p}).
\]
Verschiedene Werte von $p$ koennen als Hyperparameter verglichen werden.

\section{Evaluationsgroessen}

\subsection{Hold-out-Testfehler}
Der echte Testfehler wird auf den letzten, unberuehrten Beobachtungen berechnet:
\[
  Err_{\text{test}}
  =
  \frac{1}{T_{\text{test}}}
  \sum_{t \in \mathcal{T}_{test}}
  (y_t - \hat{y}_t)^2.
\]

\subsection{RMSE}
\[
  RMSE =
  \sqrt{
  \frac{1}{n}
  \sum_{i=1}^{n}
  (y_i - \hat{y}_i)^2
  }.
\]

\subsection{Bias der Fehlerschaetzung}
Fuer ein CV-Verfahren ist der Bias die Differenz zwischen geschaetztem
CV-Fehler und echtem Hold-out-Testfehler:
\[
  Bias =
  \frac{1}{R}
  \sum_{r=1}^{R}
  \left(
  \widehat{Err}_{CV,r} - Err_{\text{test},r}
  \right).
\]

\subsection{Varianz der Fehlerschaetzung}
\[
  Var(\widehat{Err}_{CV})
  =
  \frac{1}{R-1}
  \sum_{r=1}^{R}
  \left(
  \widehat{Err}_{CV,r}
  - \overline{\widehat{Err}_{CV}}
  \right)^2.
\]

\subsection{Model-Selection Accuracy}
Die Model-Selection Accuracy misst, wie oft ein CV-Verfahren das Modell
auswaehlt, das auf dem Hold-out-Testset tatsaechlich am besten ist:
\[
  Accuracy =
  \frac{1}{R}
  \sum_{r=1}^{R}
  \mathbb{I}
  \left(
  \hat{m}_{CV,r} = m^{*}_{test,r}
  \right).
\]

\section{Genauer Arbeitsablauf}

\subsection{Phase 1: Theorie und Eingrenzung}

\begin{enumerate}
  \item Vorlesungsfolien zu Resampling Methods und Model Selection sichten.
  \item Definitionen fuer k-fold, LOOCV, Rolling-Origin, Blocked CV und h-block CV sauber formulieren.
  \item Begruenden, warum Standard-CV bei Zeitreihen problematisch sein kann.
  \item Forschungsfrage und Vergleichskriterien final festlegen.
\end{enumerate}

\subsection{Phase 2: Simulationsdesign}

\begin{enumerate}
  \item Daten-generierende Prozesse festlegen:
  AR, MA, ARMA, Trend, Saisonalitaet.
  \item Parameter festlegen, z.B. $T = 250$,
  $T_{\text{train}} = 175$, $T_{\text{test}} = 75$.
  \item Anzahl der Monte-Carlo-Wiederholungen festlegen:
  zunaechst kleiner Testlauf, z.B. $R = 20$,
  danach Hauptsimulation, z.B. $R = 500$.
  \item Kandidaten-Lags definieren, z.B. $p \in \{1,2,5,10\}$.
  \item h-Werte fuer h-block CV definieren, z.B. $h \in \{1,5,10\}$.
\end{enumerate}

\subsection{Phase 3: Implementierung}

\begin{enumerate}
  \item Funktionen fuer die Simulation der DGPs schreiben.
  \item Funktion fuer die Erstellung von Lag-Features schreiben.
  \item CV-Splitter implementieren:
  k-fold, LOOCV, Rolling-Origin, Blocked CV, h-block CV.
  \item Modellfunktionen implementieren:
  lineare Regression, Ridge, Lasso, tree-based Modell.
  \item Evaluationsfunktionen implementieren:
  RMSE, Bias, Varianz, Model-Selection Accuracy, Laufzeit.
\end{enumerate}

\subsection{Phase 4: Smoke-Test}

\begin{enumerate}
  \item Einen sehr kleinen Testlauf ausfuehren, z.B. ein DGP, zwei Modelle,
  zwei CV-Verfahren und $R = 5$.
  \item Pruefen, ob alle Splits zeitlich korrekt sind.
  \item Pruefen, ob keine Daten aus der Zukunft ins Training gelangen.
  \item Erste Ergebnisdatei speichern und kontrollieren.
\end{enumerate}

\subsection{Phase 5: Hauptsimulation}

\begin{enumerate}
  \item Alle DGPs und Parameterkombinationen durchlaufen.
  \item Pro Simulation:
  \begin{enumerate}
    \item Zeitreihe simulieren.
    \item Train-Test-Split erstellen.
    \item Fuer jedes CV-Verfahren und jedes Modell CV-Fehler berechnen.
    \item Je CV-Verfahren das Modell mit kleinstem CV-Fehler auswaehlen.
    \item Ausgewaehltes Modell auf dem Hold-out-Testset evaluieren.
    \item CV-Fehler, Testfehler, gewaehltes Modell und Laufzeit speichern.
  \end{enumerate}
  \item Ergebnisse in einer strukturierten Tabelle zusammenfuehren.
\end{enumerate}

\subsection{Phase 6: Auswertung und Visualisierung}

\begin{enumerate}
  \item Bias der CV-Fehler je Verfahren berechnen.
  \item Varianz der CV-Fehler je Verfahren berechnen.
  \item Test-RMSE je Verfahren vergleichen.
  \item Model-Selection Accuracy je Verfahren auswerten.
  \item Laufzeiten vergleichen.
  \item Visualisierungen erstellen:
  Boxplots, Balkendiagramme, Tabellen und ggf. Heatmaps.
\end{enumerate}

\subsection{Phase 7: Interpretation}

\begin{enumerate}
  \item Pruefen, wann klassische k-fold CV zu optimistisch ist.
  \item Pruefen, ob Rolling-Origin stabilere Fehlerschaetzungen liefert.
  \item Pruefen, ob h-block CV den Bias reduziert, aber eventuell Varianz erhoeht.
  \item Trade-off zwischen Bias, Varianz und Laufzeit diskutieren.
  \item Praktische Empfehlung fuer Forecasting-Probleme formulieren.
\end{enumerate}

\section{Geplanter Output}

\begin{itemize}
  \item Reproduzierbarer Code fuer Simulation und Evaluation.
  \item Tabellen mit Bias, Varianz, RMSE, Accuracy und Laufzeit.
  \item Grafiken fuer Hausarbeit und Praesentation.
  \item Hausarbeit mit Methodik, Ergebnissen und Fazit.
  \item Praesentation mit klarer Empfehlung fuer die Anwendung.
\end{itemize}

\section{Kurze Erwartung}

Es ist zu erwarten, dass zufaellige k-fold CV und LOOCV bei stark abhaengigen
Zeitreihen zu optimistische Fehlerschaetzungen liefern koennen. Verfahren, die
die zeitliche Ordnung explizit beruecksichtigen, insbesondere Rolling-Origin,
Blocked CV und h-block CV, sollten realistischere Schaetzungen liefern. Welches
Verfahren insgesamt am besten abschneidet, haengt vermutlich von der Staerke der
Autokorrelation, Trend/Saisonalitaet, Stichprobengroesse und Modellklasse ab.

\end{document}
